\honest{A Technical Explanation of Technical Explanation}{Eliezer Yudkowsky}{2005}

This essay follows my ``Intuitive Explanation'' of Bayes's theorem. That one avoided controversy; this one ``willfully walks into the whirling helicopter blades,'' and I warn you that it does not represent ``the unanimous consensus of Earth's entire planetary community of Bayesian researchers.'' \nb{The essay is dated 2005, when it first appeared on Yudkowsky's website. The text is the later book revision, which cites posts from 2007 and 2008 and sources dated as late as 2011.}

My question is: ``What is the difference between a technical understanding and a verbal understanding?'' As a child, from popular physics, I thought I knew that sound was waves of air, light waves of electromagnetism and matter waves of probability amplitude; only after working through the wave equation in the \textsc{Feynman Lectures} did I understand ``sound is waves'' as a physicist does. That is the difference; and what I just gave you was a verbal explanation.

So here is the technical one. Think of probability as clay, a fixed amount to spread over the outcomes, say the colours a light might flash; you cannot give red 75 per cent and blue 80. To make you spread it honestly, I pay you according to how much you bet on the colour that won. Paying you the bet itself is a bad rule: if red comes up six times in ten and blue and green twice each, honest bets of 60, 20 and 20 cents earn 44 cents a round on average, but a whole dollar on red earns 60. A proper rule is one under which your best bet is your true probability. One such rule is to pay ``a dollar minus the squared error of the bet.'' \nb{For a light with three colours, as here, that rule is not proper: with frequencies of 30, 20 and 50 per cent, bets of about 35, 3 and 61 cents earn more on average than honest ones. The essay's footnote checks only two colours. Brier's rule of 1950, which charges the error on every colour, is proper.} I set it aside anyway, because ``Ordinary statisticians take the squared error of everything in sight, but not Bayesian statisticians.''

What I require instead is that your total score be the same whether two presses of the button are scored separately or together. Since probabilities multiply, ``The only possible scoring rule is'' the logarithm of the probability you gave to what happened. The base is arbitrary: in decibels, a probability of 0.01 scores $-20$; in bits, 0.25 scores $-2$. Scores are negative, which is fine; imagine the experimenter pays you a hundred dollars plus your score. With probabilities of 25, 50 and 25 per cent, your expected score is $-1.5$ bits. \nb{The derivation holds. The logarithmic rule was proposed by I.~J. Good in 1952.} It rewards calibration. Someone who says ``98\% certain'' about canola oil usually means a feeling; a calibrated 98 per cent means being wrong about twice in a hundred similar questions, and someone who says 98 per cent a thousand times and is wrong only five times should have said 99.5. The rule rewards accurate calibration, ``encouraging neither humility nor arrogance.'' It also rewards discrimination, telling right answers from wrong. Saying 50 per cent to everything is perfectly calibrated and ``the cheater's way out.'' On ten yes-or-no questions it gives the whole outcome a probability of 1 in 1,024; saying 90 per cent and being wrong twice gives 0.4 per cent, worse calibrated but better; saying 80 per cent with the same answers gives 0.6 per cent; and 99 per cent on ten right answers gives about 90 per cent. Poor calibration can be fixed by a simple transform, turning ``million-to-one'' into ``nine-to-one,'' but no transform improves discrimination; as Yates and colleagues wrote, good discrimination ``demands access to solid, predictive evidence.'' So rationality is not humble confession of helplessness: confess ignorance when ignorant and confidence when confident. A fair coin assigns every sequence of twenty flips about a millionth, $-60$ decibels, and that does not falsify it, since no other theory does better; but if someone had sealed that exact sequence in an envelope beforehand at 99 per cent, I would suspect the coin, provided she sealed only one envelope. \nb{Splitting a score into these parts is Murphy's decomposition of 1973, not cited.}

Next, a counter that shows a number from 0 to 99, perhaps the last two digits of the price of walnuts. A precise theory bets 90 per cent on 51; a vague one bets 90 per cent on the fifties, 9 per cent on each. When 51 comes up, the precise theory gains odds of ten to one, from 50 to about 91 per cent. The vague theory loses because it is timid; theories should be bold. No vague theory can claim every result from 50 to 59 as strongly, which would need 900 per cent, provided predictions are recorded in advance. If my priors are calibrated, Bayes keeps my posteriors calibrated. But if precise and vague hypotheses are equally likely as classes, there are a hundred precise ones and ten vague, so the prior odds are 1 to 10 for ``51'' against ``fifties,'' and one 51 brings them only to even, as common sense expects. A coin that shows HHTTH:TTTTH may be fixed, but a 1 per cent prior of fixedness covers every sequence, so a random-looking result tells you only which sequence it would be fixed to; a second run of the same sequence would make fixedness ten to one. That is why Bayesians need priors. Along the way: classical statistics ``insists on paying attention only to likelihoods'' and so always overfits. \nb{The usual complaint against classical statistics is the reverse, that it violates the likelihood principle; and classical methods such as Akaike's criterion charge a hypothesis for each parameter fitted to the data.} \nb{These figures, and the essay's other worked numbers, check to rounding.} A theory of complete ignorance, spreading probability evenly, loses to the vague theory on 51; with prior odds of 1 to 10 to 200 for the precise, vague and ignorant theories, two 51s leave them at 89, 9 and 2 per cent. A stupid theory that bets 90 per cent on 0 to 9 is falsified at once. A theory that fits every outcome is the same as ignorance, and a theory that bets on the wrong range does worse still: ``Ignorance is better than anti-knowledge.'' (So an AI algorithm that improves when noise is added must be doing something worse than random.) On the results 52, 51 and 58, the vague theory scores about $-30$ decibels, the precise and ignorant $-60$ and the stupid $-90$; with priors the vague still wins. Even the best theory scores low; theories are approximations.

Humans are not like this. They ``bizarrely attempt to defend hypotheses,'' which ``has no analogue in the laws of probability theory or decision theory'': you may test ideas you are unsure of, but not prefer an outcome. If A is evidence for B, not-A must be evidence against it; yet people want every result to prove their theory, like Spee's witch-hunters, who would ``feel disgraced if it acquitted a woman.'' \nb{Lakatos argued that defending a theory's core while adjusting auxiliary hypotheses is how good science proceeds. His standard example is the Newtonian search for an unseen planet, which this essay later tells as Newton's triumph.} Humans judge by ``fit'', which is not conserved, and they judge after the fact; a psychoanalyst can make any behaviour fit. The students who explained why the side of a plate near a radiator was cooler lost their chance at bewilderment; equations of diffusion would have given a sharp prediction. You can now calculate exactly how technical an explanation is, though for vague ones only a superhuman mind could do it. So everyone needs at least one technical subject, or they may think, like Heraclitus, that ``All is Fire'' explains something; I would teach Bayesian probability in high school. In technical words: ``Verbal explanation runs on psychological evaluation of unconserved post facto compatibility instead of conserved ante facto probability density.'' Then I admit that those words teach nothing to someone who cannot do the math, say ``Bayes'' nine times, and add that ``if you know the math, you can be as silly as you like, and still technical'': a scientist should be able to make discoveries in a clown suit or lecture on helium. A useful model is knowledge you can compute in reasonable time: a model of a 747 that slightly violates conservation of momentum may predict more cheaply, without the plane violating it, and atomic theory, far from competing with aerodynamics, mandates it as an approximation.

Now I lower my own standard. Many models in evolutionary theory are quantitative, such as the drift of non-functional changes allowed by the redundant genetic code, where 64 codons code for 20 amino acids and a stop. But in the nineteenth century, before anyone knew of DNA or that humans and chimpanzees share most of their genes, scientists still flocked to natural selection. Physics confirms General Relativity to one part in $10^{14}$; can evolution match that? Dennett called Darwin's the single greatest idea. Yet when precise physics said the Sun could have burned only three thousand years on chemical energy or 40 million on gravitational, too short for evolution, and Kelvin wrote that this seemed ``sufficient to disprove'' descent by natural selection, nineteenth-century evolution, which I describe as ``without a scrap of quantitative modeling,'' turned out right. \nb{The first edition of \textsc{On the Origin of Species} estimated three hundred million years for the erosion of the Weald, a number bearing on exactly this question.} ``History records who won.'' Predictions right 80 per cent of the time can build a huge lead over ignorance. Reality is consistent, so both mountains of evidence were right in their domains; nineteenth-century physics did not know about nuclear reactions. So vague, ``semitechnical'' theories can be well confirmed, and when two well-confirmed theories clash, one is being misapplied. But telling good semitechnical theories from confused ones takes skill, which is why we have Science. \nb{Huxley answered Kelvin along these lines in 1869. In Wikipedia's summary, he suggested that Kelvin's calculations ``appeared precise in themselves but were based on faulty assumptions.''} Evolution also explained facts already known, and that was ``correct as probability theory.'' The rule of advance prediction is ``a moral custom and not a theorem,'' kept to stop people tinkering after the fact: ``The social process of science is a set of legal conventions to keep people from cheating on the math.'' Semitechnical theories have a cost: Darwin made remarkably few errors, but his successors, bright enough only to accept the theory, spoke as late as the 1960s of evolution working for the good of the species. Technical theories are far better, but Nature does not always allow them at once, and scientific controversy is mostly about semitechnical theories or nonsense posing as them. So prefer textbooks to news: ``gravity results from the curvature of spacetime'' is fascinating words without the math. Follow a controversy only if it is your field or affects your life, and do not apply stricter standards to theories you dislike, or every new flaw you learn to spot makes you stupider. \nb{Mill and Keynes held the same view; Whewell and, in Bayesian terms, Maher held the opposite. Standard Bayesian confirmation theory also has trouble crediting old evidence at all, the ``problem of old evidence'' (Glymour, 1980). The essay mentions neither dispute.}

A classic sign of a poor hypothesis is its effort to avoid falsification, like Sagan's dragon in the garage; but the claimant knows in advance which excuses will be needed, so anticipates exactly what the skeptic does. I ask which experiences I would anticipate, not which facts I would believe. The flip side: I then tell of an argument with a friend who said punctuated equilibrium showed something wrong with evolutionary theory. I found \textsc{Homo heidelbergensis} as an intermediate, but my friend said the changes were still too sharp, and when I suggested that selection pressures vary, said I was making excuses in advance. Yet if the record had been smooth, might he not have said a noisy process could not produce that? Critics must state exactly what the standard theory predicts badly, and how they know. Ask a theory's advocates for its predictions, in advance, and let detractors cross-examine. Models include noise: a precise theory that puts 90 per cent on 51 is not refuted by one 82 among nine 51s; that is an advance prediction, not an excuse. As for me, ``Maybe there are models in the evolutionary family that would make advance predictions of steadiness or variability, but if so, I don't know about them.'' I admit that ``I rendered even that verdict after seeing the data.'' \nb{Eldredge and Gould had derived punctuated equilibrium in 1972 from Mayr's standard account of speciation, which supports Yudkowsky's side of the argument.}

Mercury's orbit defied Newton, who predicted 5,557 seconds of arc a century against 5,600 observed, and nobody could know in advance that it would take Einstein to explain it. Uranus defied Newton too, and the answer was Neptune, found within a degree of where Le Verrier said. Neptune was ``the first observation that confirmed an advance prediction of Newtonian gravitation.'' \nb{Halley's comet, predicted by Halley with Newton's laws and seen again in December 1758, is the usual earlier case.}

Popper erred ``in thinking that falsification was qualitatively different from confirmation,'' though he saw a real quantitative difference. He was impressed that Einstein's theory was risky while the theories of Freud and Adler could explain any behaviour and so predicted nothing. I quote an encyclopedia on his view and translate it: ``the more outcomes a model prohibits, the more probability density the model concentrates in the remaining, permitted outcomes.'' \nb{The same encyclopedia entry reports that Popper held ``the more improbable a theory is the better it is scientifically,'' the opposite of the Bayesian preference for the theory with the higher posterior. The notes on ``An Intuitive Explanation'' give Popper's view of confirmation.} Phlogiston and vitalism explained everything, so they were ignorance in disguise. But beware checklists: a fundamental element is not a mysterious one. Physics takes the electromagnetic field as fundamental, yet governs it by simple computable rules, while crackpots demand an ``underlying mechanism'' and offer only words. Fundamentals should be simple: oxygen is a good one, life a bad one, since psychologically simple things like flame and moving flesh hide vast complexity.

Then a parable. Believers in St. Peter's judgement get no computable scoring rule. In the religion of ``Bayesianity'', Laplace multiplies every probability you ever assigned, and that product is your eternal score. Say 99.9 per cent each morning that the Sun will rise, and your score falls by a factor of 0.7 a year; at 99.999 per cent, after seventy years your score is still 77.4 per cent of the best possible. Assigning probability 1 to anything is its ``sole mortal sin'', since one miss makes the product zero, and you spend eternity beside a believer in flying saucers. If an AI might take the Sun apart some night in the next ten years, with even odds, the right nightly figure is about 99.98 per cent; then say exactly that. \nb{Dennis Lindley called this Cromwell's rule.} You cannot game the score by modesty or boldness, only by guessing better: ask not which excuses you have but which you expect to need. I admit that you could also game it if ``you could commit suicide when you turned five,'' and I drop the religion there.

Advance prediction matters because people find it hard to obey probability after seeing the data. Darwin saw the similarity of species before thinking of natural selection, and Newton's laws explained orbits and tides already observed; Neptune came long after the theory was accepted. A numerologist cannot predict next week's Mega Ball, but explains why last week's seven was inevitable; asked to bet, he cannot give more than 1 in 52 to seven without taking it from others, and a listener grows suspicious as he explains why each of balls one to twelve fits.

Then imagine waking with a blue tentacle for an arm. I would not explain it; it will not happen. Divine intervention, aliens or hallucination ``fit'' anything and equal ignorance; why that tentacle, that morning? And if aliens did it, what will they do tomorrow? A good explanation is one that would make me nervous now, before it happens, and explaining events you do not expect ``necessarily violates the laws of probability theory.'' In short, ``to explain is to anticipate.'' A witch who transports me into a webcomic would predict the tentacle, but its prior is effectively zero, so it explains nothing. With calibrated priors, among two million beings who meet such strangeness, ten owe it to improbable hypotheses (prior a thousandth, likelihood a hundredth) and a thousand to probable ones (a hundredth and a tenth), so a real tentacle would come from something more normal than webcomic witchery that I did not see coming. ``Reality dishes out experiences using probability, not plausibility.'' \nb{Asking which hypothesis would best explain an unlikely event is an ordinary conditional question, and a few paragraphs later the essay answers it by counting two million beings across the universe. What probability rules out is holding an explanation probable while not expecting what it predicts.}

Last, a verse, ``Since the beginning / Not one unusual thing / Has ever happened,'' and a request to ``take joy in the merely real.''
